\scanpage{26}
\mspage{13}

à prouver que $\widetilde u_\alpha\in C$. Lorsque on suppose de plus $u_\bullet\in W_{I^\circ}$ i.e les $u_i\in W$, il faut de plus prouver que $\widetilde u_\alpha\in TC=C\cap W$. La prop.~2 et son corollaire 1 résulteront de là. Pour prouver $\widetilde u_\alpha\in C$ \deleted{\unsure{i.e.}} resp. $\widetilde u_\alpha\in TC$, on \deleted{utilise sa \unsure{proc.}} peut supposer $I_\alpha=I$ (quitte à changer de notation). On pose donc

\[
\begin{aligned}
 I_\alpha&=I,\ I_{\alpha-1}=I'\\
 I_{/s_{\alpha-1}}&=I_0,\ I_0\cap I'=I'_0,
\end{aligned}
\]

de sorte qu’on a un diagramme commutatif d’inclusions ouvertes

\diagram{\begin{tikzcd}[ampersand replacement=\&]
 I_0\cap I'=I'_0\arrow[r,hook]\arrow[d,hook]\&I_0\arrow[d,hook]\\
 I'\arrow[r,hook]\& I=I_0\cup I'\, .
\end{tikzcd}}

On utilisera seulement le fait que $I_0$, donc $I'_0$, est fini, i.e que \inserted{dans la \unsure{config.|cond.} ($I'\subset I$),} $I$ peut s’écrire sous la forme $I=I'\cup I_0$, avec $I_0$ un ouvert \uline{fini}, pour en conclure que dans le diagramme

\diagram{\begin{tikzcd}[ampersand replacement=\&,column sep=large]
 X_{I'}\arrow[r,"{\tau^X_{I,I'}}"]\arrow[d,"u_{I'}"']\arrow[dr,phantom,"\text{cocart}" description]
 \&\deleted{\unsure{I}}X_{\deleted{\unsure{1|{}}}I}\arrow[d,"{\unsure{v|u}_{I',I}}"]\arrow[dd,bend left=42,"u_I"]\\
 Y_{I'}\arrow[r]\arrow[dr,"{\tau^{Y}_{I,I'}}"']\&\widetilde Y_I\arrow[d,"{\widetilde u=u_{I',I}}"]\\
 \&Y_I
\end{tikzcd}}

on a $\widetilde u_{\deleted{\unsure{I}}}\in C$ resp $\widetilde u\in TC=C\cap W$.
% archive sheet number: 25

